\scanpage{61}
\deleted{\unsure{4}}
% archive foliation: 63
b) Recollement de $\overline H$ à lui $=$ suivant
un isom. rectifiable $[0,1]\xrightarrow{\sigma}[-1,0]$ \inserted{$(\sigma(0)=0)$},
regardé au voisinage du pt 0.

\margin{\uline{12.1}\quad\deleted{\unsure{a)}} c) Soit $C$ une
courbe rectifiable dans \deleted{\unsure{un voisinage de}}\inserted{un ouvert} $U$,
\deleted{\unsure{image de}} $\unsure{\mathbb C}$,
toute \deleted{\unsure{application}} fonction \inserted{complexe}
continue sur $U$ qui est
\unsure{continue|analytique} dans $U\setminus C$, \unsure{est en}
fait holomorphe dans $U$.}

\textcircled{4} \uline{Recollements \deleted{bord à bord}\inserted{des disques}}.
\margin{Le recollement ``bord à bord''}
Cela correspond,
dans le th. 4, au cas où $K$ est une
variété \inserted{(compacte)} de dim 1, $S=\varnothing$. \deleted{Il} Pour avoir
``des \unsure{esp.} de modules'', il est naturel
de dégager des conditions restrictives
naturelles sur l'involution $\sigma$ de
$\delta\widetilde X$.

Pour ceci, \deleted{\unsure{rapp}} on se rappelle que
pour une surface conforme à bord $X$
\inserted{(au lieu de $\widetilde X$)}, et une comp. connexe du bord, soit $C$,
il revient au $=$ de se donner une
structure de droite projective \inserted{rectifiable} sur $C$,
ou un ``bouchage de trou'' en $C$
par un disque \inserted{de bord $C$}, les propriétés de régularité
(\unsure{p.r. à} la structure analytique réelle de $C$)
de la structure de droite projective, reflétant
